\documentclass[12pt]{article}
\usepackage[utf8]{inputenc}
\usepackage[T1]{fontenc}
\usepackage{hyperref}
\title{Ulumanda' Language Resources for K-12 Curriculum Development}
\author{Compiled by Claude Sonnet 5 (claude-sonnet-5) for Andr\'as Kornai}
\date{}
\begin{document}
\maketitle

\section{Introduction}
Ulumanda' (ISO 639-3 ulm) is an Austronesian language of West Sulawesi, Indonesia, in the Pitu Ulunna Salu group. Speakers are mostly in Majene Regency, with some in Polewali Mandar and Mamuju. About 30,000 speakers are cited from 1986. Dialects named include Botteng, Sondoang and Tapalang. The searches found no dictionary, primer or school materials.

\section{Primary Educational Resources}
None found.

\section{Linguistic and Reference Materials}
\begin{itemize}
\item Wikipedia, Ulumanda' language: \url{https://en.wikipedia.org/wiki/Ulumanda%CA%BC_language}
\item Brodkin (2019), comparative paper on West Sulawesi languages with Ulumanda' data and basic clause structure (UC Santa Cruz; not a grammar of Ulumanda' alone): \url{https://people.ucsc.edu/~ddbrodki/PDFs/Brodkin_NUS_19_1.pdf}
\item Omniglot, Ulumanda' writing page: \url{https://www.omniglot.com/writing/ulumanda.htm}
\item Omniglot, Ulumanda' numbers: \url{https://omniglot.com/language/numbers/ulumanda.htm}
\end{itemize}
The search summary mentioned Zobel (2018) on the vowel inventory, but no URL for it was seen.

\section{Religious and Community Resources}
\begin{itemize}
\item Joshua Project, Ulumanda in Indonesia: \url{https://joshuaproject.net/people_groups/15705}
\item Praying for Indonesia (profile page): \url{https://www.prayingforindonesia.com/?p=1184}
\end{itemize}
No scripture translations were confirmed.

\section{Digital and Audio Resources}
None found. The Endangered Languages Alliance page for Mamuju concerns a different language and is not Ulumanda' material.

\section{Curriculum Development Context}
Ulumanda' is still spoken by all adults, so a community-led curriculum is realistic. The found materials are thin: scholarly papers, a number chart and an alphabet chart, and no dictionary or textbook. Bahasa Indonesia dominates schooling in the area. Initial work would need documentation with local speakers, such as a word list, an agreed orthography and recorded stories.

\end{document}